% Generated by scripts/assemble.py; corrections belong in source/corrections.json.
\examhead{150}{国防科技大学 2024 年数学分析与高等代数试卷}{143}
% p144-b01; source PDF page 144
\par\Needspace{5\baselineskip}
\blockanchor{p144-b01}
\majorquestion{一}证明题. 每题 15 分, 共 120 分.\par
% p144-b02; source PDF page 144
\par\Needspace{6\baselineskip}
\blockanchor{p144-b02}
\question{1}给定两正数 $a_1, b_1, a_1 > b_1$. 作出其等差中项 $a_2 = \frac{a_1 + b_1}{2}$ 与等比中项 $b_2 = \sqrt{a_1b_1}$. 一般地, 定义\par
\nopagebreak[4]
% p144-b03; source PDF page 144
\blockanchor{p144-b03}
\displayformula{a _ {n + 1} = \frac {a _ {n} + b _ {n}}{2}, b _ {n + 1} = \sqrt {a _ {n} b _ {n}}, n = 1, 2, \dots}
% p144-b04; source PDF page 144
\blockanchor{p144-b04}
证明： $\lim_{n\to \infty}a_n$ 与 $\lim_{n\to \infty}b_n$ 存在且相等.\par
% p144-b05; source PDF page 144
\blockanchor{p144-b05}
\question{2}设 $\sum_{n=1}^{\infty} a_n$ 为发散的非负项级数, $S_n = \sum_{k=1}^{n} a_k$ 且 $a_1 > 0$. 证明: 级数 $\sum_{n=1}^{\infty} \frac{a_n}{S_n^2}$ 收敛.\par
% p144-b06; source PDF page 144
\blockanchor{p144-b06}
\question{3}证明：当 $0 < \varepsilon < 1$ 时，方程 $y = x - \varepsilon \sin x$ 可以唯一确定一个函数 $x = f(y)$，它是定义在 $(- \infty, +\infty)$ 上的严格单调递增函数。\par
% p144-b07; source PDF page 144
\blockanchor{p144-b07}
\question{4}证明: 当 $a > 0$ 时, 含参量反常积分 $\int_{0}^{+\infty} x e^{-xy} \mathrm{d}y$ 在 $[a, +\infty)$ 上一致收敛, 但在 $(0, +\infty)$ 上不一致收敛.\par
% p144-b08; source PDF page 144
\blockanchor{p144-b08}
\question{5}设函数 $f(x)$ 在 (0,1] 内连续且可导，并且 $\lim_{x\to 0^+}\sqrt{x} f'(x) = a$. 证明：$f(x)$ 在 (0,1] 内一致连续.\par
% p144-b09; source PDF page 144
\blockanchor{p144-b09}
\question{6}设 $n$ 阶实方阵 $A$ 满足 $A^3 = A$. 证明: 若对于任意实向量 $\alpha$, 均有 $\alpha^T A^T A \alpha \leq \alpha^T \alpha$, 则 $A$ 是实对称矩阵.\par
% p144-b10; source PDF page 144
\blockanchor{p144-b10}
\question{7}设 $A, B$ 为 $n$ 阶矩阵. 证明:\par
% p144-b11; source PDF page 144
\blockanchor{p144-b11}
(1) 若 $AB = O$, 则 $r(A) + r(B) \leq n$.\par
% p144-b12; source PDF page 144
\blockanchor{p144-b12}
(2) 若 $A^2 = E$, 则 $r(A + E) + r(A - E) = n$.\par
% p144-b13; source PDF page 144
\blockanchor{p144-b13}
\question{8}设 $A$ 为复方阵, $p(x)$ 为 $E - \overline{A} A$ 的特征多项式, 其中 $\overline{A}$ 是 $A$ 的共轭矩阵, 证明: $p(x)$ 必为实系数多项式.\par
% p144-b14; source PDF page 144
\par\Needspace{5\baselineskip}
\blockanchor{p144-b14}
\majorquestion{二}解答题. 每题 15 分, 共 30 分.\par
% p144-b15; source PDF page 144
\blockanchor{p144-b15}
\question{1}一质点在变力 $F = (xy + e^{-x}\cos y, x^2 + e^{-x}\sin y)$ 的作用下，从原点沿抛物线 $y = x^2$ 运动到点 $A(1,1)$，再沿直线运动到点 $B(0,2)$，求该过程中变力所做的功.\par
% p144-b16; source PDF page 144
\par\Needspace{6\baselineskip}
\blockanchor{p144-b16}
\question{2}当 $a, b$ 满足何条件时, 方程组有解? 有解的情况下, 求出全部解.\par
\nopagebreak[4]
% p144-b17; source PDF page 144
\blockanchor{p144-b17}
\displayformula{\left\{ \begin{array}{l} a x _ {1} + x _ {2} + x _ {3} = 4; \\ x _ {1} + b x _ {2} + x _ {3} = 3; \\ x _ {1} + 2 b x _ {2} + x _ {3} = 4. \end{array} \right.}
